\documentclass[12pt, a4paper]{article}

\usepackage[utf8]{inputenc}
\usepackage[T1]{fontenc}
\usepackage[brazil]{babel}
\usepackage{amsmath, amssymb}
\usepackage{graphicx}
\usepackage{float}
\usepackage{booktabs}
\usepackage{geometry}
\usepackage{setspace}
\usepackage{xcolor}
\usepackage{fancyhdr}
\usepackage{hyperref}
\usepackage{caption}
\usepackage{indentfirst}

\geometry{left=2.5cm, right=2.5cm, top=2.5cm, bottom=2.5cm}
\setstretch{1.3}
\setlength{\parindent}{1.5em}
\setlength{\parskip}{0.4em}

\hypersetup{
    colorlinks=true,
    linkcolor=blue!70!black,
    citecolor=blue!70!black,
    urlcolor=blue!70!black
}

\pagestyle{fancy}
\fancyhf{}
\setlength{\headheight}{14.5pt}
\addtolength{\topmargin}{-2.5pt}
\rhead{Projeto 2 --- ARMA-GARCH PETR4}
\lhead{Séries Temporais}
\cfoot{\thepage}

\begin{document}

% ============================================================
% CAPA
% ============================================================
\begin{titlepage}
\centering
\vspace*{2cm}
{\Large UNIVERSIDADE ESTADUAL DE CAMPINAS}\\[0.2cm]
{\large Séries Temporais --- ME607}\\[0.2cm]
{\LARGE Modelagem da Volatilidade e Análise de Intervenção em Retornos da Petrobras (PETR4):\\[0.4cm] Impactos de Eventos Político-Econômicos no Período 2010--2026}\\[2cm]
{\Large Projeto 2 --- Análise de Dados}\\[3cm]
{\large
Nome: Arthur Lourenço Narcizo da Silva\\[0.2cm]
}
\vfill
{\large Campinas, \today}
\end{titlepage}

\tableofcontents
\newpage

% ============================================================
% RESUMO
% ============================================================
\section*{Resumo}
\addcontentsline{toc}{section}{Resumo}

Modelei os log-retornos diários da PETR4 entre 04/01/2010 e 29/04/2026, num total de 4053 observações, usando a família ARMA-GARCH. A análise descritiva mostrou caudas pesadas (curtose em excesso de 13,35), assimetria negativa de $-0{,}93$ e forte heterocedasticidade condicional, com o teste ARCH-LM rejeitando homocedasticidade a $p$-valor inferior a $10^{-130}$. Estimei três modelos por máxima verossimilhança: GARCH(1,1)-$t$, EGARCH(1,1)-$t$ e EGARCH(1,1)-$t$ com regressores de intervenção para eventos político-econômicos.

O modelo final, escolhido pelo menor AIC, capturou efeito alavancagem significativo ($\gamma = -0{,}030$, $p < 0{,}001$) e estimou impacto pontual de $-6{,}50$ pontos percentuais no retorno do dia da demissão de Castello Branco. Calculei o Value-at-Risk paramétrico nos níveis de 95\% e 99\% e o teste de Kupiec não rejeitou a adequação do modelo a 99\% ($p = 0{,}488$), o que permite usar o modelo para gestão de risco de mercado.

\vspace{0.5em}
\noindent Palavras-chave: séries temporais financeiras; GARCH; EGARCH; análise de intervenção; Value-at-Risk; PETR4.

\newpage

% ============================================================
% 1. INTRODUÇÃO
% ============================================================
\section{Introdução}

A modelagem de séries temporais financeiras é uma das áreas mais desenvolvidas e aplicadas da econometria moderna, em razão de sua relevância para gestão de risco, precificação de derivativos e alocação de portfólios. Séries de retornos exibem propriedades estatísticas peculiares, conhecidas na literatura como fatos estilizados, que tornam inadequada a aplicação direta de modelos clássicos baseados em homocedasticidade e normalidade \cite{morettin2017}.

A ação preferencial da Petrobras é um caso emblemático no mercado brasileiro. Por sua liquidez, peso no Ibovespa e exposição a riscos políticos, regulatórios e geopolíticos, a série de retornos apresenta clusters de volatilidade pronunciados e várias quebras estruturais associadas a eventos discretos --- crises de governança, mudanças de comando, alterações na política de preços de combustíveis. Estudar esses fenômenos com modelos rigorosos tem interesse acadêmico e prático.

O objetivo deste trabalho é construir um modelo estatístico para a média e a volatilidade condicionais dos log-retornos diários da PETR4 no período 2010--2026, capaz de capturar os fatos estilizados das séries financeiras e de quantificar o impacto de eventos político-econômicos. Como objetivos específicos, busco realizar análises descritiva, exploratória e inferencial dos retornos; identificar e estimar modelos da família GARCH com diagnósticos formais; mensurar o efeito de eventos relevantes via análise de intervenção; e calcular o Value-at-Risk diário com posterior validação por backtesting.

A escolha da PETR4 não é arbitrária. Quis um ativo brasileiro com histórico longo, episódios de stress bem documentados e relevância prática para o mercado local --- e PETR4 cumpre as três condições com folga.\footnote{Cheguei a considerar VALE3 e ITUB4, mas a riqueza dos eventos políticos sobre a Petrobras tornou a análise de intervenção mais informativa.}

O texto está organizado da seguinte forma. A Seção~\ref{sec:teoria} apresenta o referencial teórico, com ênfase nos modelos GARCH e EGARCH, na análise de intervenção e no Value-at-Risk. A Seção~\ref{sec:metodo} detalha a metodologia, os dados e o software utilizado. A Seção~\ref{sec:descritiva} traz a análise descritiva. A Seção~\ref{sec:modelo} documenta a estimação e a seleção do modelo final. A Seção~\ref{sec:interv} é dedicada à análise de intervenção. A Seção~\ref{sec:diag} apresenta os diagnósticos. A Seção~\ref{sec:var} calcula o VaR e o submete a backtesting. A Seção~\ref{sec:disc} traz a discussão crítica e a Seção~\ref{sec:conc} conclui.

% ============================================================
% 2. REFERENCIAL TEÓRICO
% ============================================================
\section{Referencial Teórico}
\label{sec:teoria}

\subsection{Fatos estilizados de séries financeiras}

A literatura empírica em finanças, discutida em Morettin~\cite{morettin2017} e em Shumway e Stoffer~\cite{shumway2017}, documenta propriedades recorrentes em séries de retornos. Entre as principais, destacam-se a ausência de autocorrelação linear nos retornos, contrastada pela autocorrelação significativa nos retornos ao quadrado --- o que indica dependência não linear. Há também evidência empírica robusta de caudas mais pesadas que as da normal (leptocurtose) e de clusters de volatilidade, em que períodos de oscilações grandes tendem a ser seguidos por mais oscilações grandes. Observa-se ainda o efeito alavancagem: assimetria entre choques positivos e negativos sobre a variância futura.

\subsection{Log-retornos}

Seja $P_t$ o preço de um ativo no instante $t$. Define-se o log-retorno em percentual como
\begin{equation}
    r_t = 100 \cdot \ln\!\left(\frac{P_t}{P_{t-1}}\right),
\end{equation}
medida preferida em finanças por sua aditividade temporal e por aproximar-se da distribuição normal sob agregação, conforme discutido em Morettin~\cite{morettin2017}.

\subsection{Modelos ARMA}

Um processo $\{r_t\}$ segue um modelo ARMA$(p,q)$ se
\begin{equation}
    \phi(B)\,r_t = \theta(B)\,a_t,
\end{equation}
onde $B$ é o operador de retardo, $\phi(B) = 1 - \phi_1 B - \cdots - \phi_p B^p$, $\theta(B) = 1 + \theta_1 B + \cdots + \theta_q B^q$ e $\{a_t\}$ é ruído branco com $\mathbb{E}(a_t) = 0$ e $\mathrm{Var}(a_t) = \sigma_a^2$. A metodologia de identificação, estimação e diagnóstico segue o roteiro clássico de Box e Jenkins, descrito em detalhe em Morettin e Toloi~\cite{morettin2018}.

\subsection{Modelos GARCH e EGARCH}

Engle~\cite{engle1982} propôs o modelo ARCH$(q)$ para descrever a heterocedasticidade condicional observada em séries financeiras --- trabalho seminal que lhe rendeu o Prêmio Nobel de Economia em 2003. Bollerslev~\cite{bollerslev1986} estendeu a formulação para o GARCH$(p,q)$:
\begin{equation}
    a_t = \sigma_t \varepsilon_t, \qquad
    \sigma_t^2 = \omega + \sum_{i=1}^{q} \alpha_i a_{t-i}^2 + \sum_{j=1}^{p} \beta_j \sigma_{t-j}^2,
\end{equation}
com $\omega > 0$, $\alpha_i, \beta_j \geq 0$ e $\sum (\alpha_i + \beta_j) < 1$ como condição de estacionariedade da variância.

Uma limitação conhecida do GARCH é tratar de forma simétrica choques positivos e negativos. Para contornar isso, Nelson~\cite{nelson1991} propôs o EGARCH$(p,q)$, definido sobre o logaritmo da variância:
\begin{equation}
    \ln \sigma_t^2 = \omega + \sum_{i=1}^{q} \big[ \alpha_i |\varepsilon_{t-i}| + \gamma_i \varepsilon_{t-i} \big] + \sum_{j=1}^{p} \beta_j \ln \sigma_{t-j}^2.
\end{equation}
Quando o coeficiente $\gamma$ é negativo e significativo, há evidência de efeito alavancagem: choques negativos elevam a volatilidade futura mais do que choques positivos de mesma magnitude. Outra vantagem do EGARCH é dispensar restrições de positividade sobre os coeficientes, já que modela diretamente o logaritmo da variância.

\subsection{Distribuição $t$-Student padronizada}

Dada a leptocurtose dos retornos financeiros, é prática estabelecida adotar inovações $\varepsilon_t \sim t_\nu$ padronizada, com variância unitária e $\nu > 2$ graus de liberdade. Quanto menor o valor de $\nu$, mais pesadas são as caudas. No limite $\nu \to \infty$, recupera-se a normal padrão.

\subsection{Análise de intervenção}

A análise de intervenção, discutida em Morettin e Toloi~\cite{morettin2020}, propõe incluir variáveis indicadoras na equação da média para modelar o impacto de eventos discretos sobre a série. As duas formas básicas de função de intervenção são o pulso e o degrau. No pulso, a indicadora $I_t$ assume valor unitário apenas no instante do evento $t = T^*$, capturando um efeito transitório de um único período. No degrau, a indicadora assume valor unitário a partir do evento --- isto é, para todo $t \geq T^*$ ---, capturando um efeito permanente sobre o nível da série. A média condicional fica então $\mu_t = c + \phi_1 r_{t-1} + \sum_k \delta_k I_{k,t}$, onde $\delta_k$ representa o efeito do $k$-ésimo evento.

\subsection{Value-at-Risk e backtesting}

O Value-at-Risk no nível $\alpha$ é o quantil-$\alpha$ da distribuição condicional de perdas:
\begin{equation}
    \mathrm{VaR}_t(\alpha) = -\big( \mu_t + q_\alpha(\nu) \cdot \sigma_t \big),
\end{equation}
onde $q_\alpha(\nu)$ é o quantil-$\alpha$ da $t$-Student padronizada \cite{morettin2017}. Para validar o modelo, usa-se o teste de cobertura incondicional de Kupiec, baseado na razão de verossimilhança
\begin{equation}
    \mathrm{LR}_{uc} = -2 \ln\!\left[ \frac{(1-\alpha)^{n-x} \alpha^x}{(1-\hat{p})^{n-x} \hat{p}^x} \right] \;\sim\; \chi^2_1,
\end{equation}
em que $x$ é o número de violações observadas, $n$ o tamanho amostral e $\hat{p} = x/n$ a taxa empírica de violações. Sob a hipótese nula de cobertura correta, a estatística segue qui-quadrado com um grau de liberdade.

% ============================================================
% 3. METODOLOGIA
% ============================================================
\section{Metodologia}
\label{sec:metodo}

\subsection{Dados}

Coletei preços de fechamento ajustados diários da PETR4 entre 04/01/2010 e 29/04/2026, totalizando 4054 pregões e 4053 log-retornos. A fonte primária foi o Yahoo Finance, e os dados foram consolidados em arquivo CSV hospedado em repositório público no GitHub: \url{https://github.com/arthurnarcizo/S-ries_Temporais_ME607}. Isso garante reprodutibilidade plena dos resultados. O preço utilizado já incorpora ajustes por proventos e desdobramentos.\footnote{Inicialmente tentei coletar os dados diretamente pela API do BCB, mas a série não cobria o período completo. O Yahoo Finance, ainda que com limitações conhecidas, foi a opção mais prática.}

\subsection{Eventos selecionados para análise de intervenção}

Após levantamento da imprensa especializada e do histórico de preços do ativo, selecionei três eventos com elevada repercussão e impacto observável (Tabela~\ref{tab:eventos}).

\begin{table}[H]
\centering
\caption{Eventos político-econômicos modelados por intervenção.}
\label{tab:eventos}
\begin{tabular}{lcl}
\toprule
Evento & Data & Tipo de dummy \\
\midrule
Demissão de Roberto Castello Branco        & 19/02/2021 & Pulso \\
Mudança da política de preços (gov.\ Lula) & 08/03/2023 & Degrau \\
Demissão de Jean Paul Prates               & 14/05/2024 & Pulso \\
\bottomrule
\end{tabular}
\end{table}

\subsection{Software e reprodutibilidade}

A análise foi conduzida em Python 3.11 com as bibliotecas \texttt{pandas}, \texttt{numpy}, \texttt{scipy}, \texttt{statsmodels} e \texttt{arch}~\cite{sheppard2024}. O código-fonte é fornecido como arquivo suplementar e a semente aleatória foi fixada em 42 para garantir resultados idênticos em execuções futuras. O fluxograma do método é apresentado em arquivo suplementar e segue o roteiro Box-Jenkins adaptado à modelagem da volatilidade condicional, com as etapas de identificação, estimação, diagnóstico, previsão (VaR) e backtesting.

% ============================================================
% 4. ANÁLISE DESCRITIVA E EXPLORATÓRIA
% ============================================================
\section{Análise Descritiva e Exploratória}
\label{sec:descritiva}

\subsection{Estatísticas resumo}

A Tabela~\ref{tab:descritivas} apresenta as estatísticas descritivas dos log-retornos. A média é próxima de zero ($0{,}0447\%$ ao dia) e o desvio-padrão é alto, $2{,}74\%$, equivalente a cerca de $43{,}4\%$ de volatilidade anualizada\footnote{Calculada como $\sigma_{\text{anual}} = \sigma_{\text{diária}} \cdot \sqrt{252} \approx 2{,}74 \cdot 15{,}87$.}. A assimetria negativa de $-0{,}93$ indica preponderância de retornos extremos no lado das perdas, e a curtose em excesso de $13{,}35$ revela caudas muito mais pesadas que as da normal. Esses padrões batem com os fatos estilizados da literatura \cite{morettin2017} e justificam tanto o uso de inovações $t$-Student quanto a adoção de modelos com volatilidade variável no tempo.

\begin{table}[H]
\centering
\caption{Estatísticas descritivas dos log-retornos diários (\%) de PETR4, 2010--2026.}
\label{tab:descritivas}
\begin{tabular}{lr}
\toprule
Estatística & Valor \\
\midrule
N (observações)    & 4053 \\
Média              & $0{,}0447$ \\
Mediana            & $0{,}0662$ \\
Desvio-padrão      & $2{,}7351$ \\
Mínimo             & $-35{,}2367$ \\
Máximo             & $20{,}0671$ \\
Assimetria         & $-0{,}9272$ \\
Curtose (excesso)  & $13{,}3484$ \\
\bottomrule
\end{tabular}
\end{table}

\subsection{Visualização da série}

A Figura~\ref{fig:serie} mostra a série de preços e dos log-retornos com os eventos modelados destacados por linhas tracejadas. No painel superior, fica visível a forte queda do preço entre 2014 e 2016 --- período do desenrolar da Operação Lava Jato e do ciclo de baixa do petróleo. No painel inferior, percebem-se nitidamente os clusters de volatilidade. Chamam atenção o período da pandemia da COVID-19 em março de 2020, quando se registrou o retorno mínimo da amostra ($-35{,}24\%$), e o salto negativo associado à demissão de Castello Branco em fevereiro de 2021.

\begin{figure}[H]
\centering
\includegraphics[width=\linewidth]{01_serie_eventos.png}
\caption{PETR4: preço de fechamento ajustado (acima) e log-retornos diários em \% (abaixo). As linhas tracejadas vermelhas indicam os eventos modelados.}
\label{fig:serie}
\end{figure}

\subsection{Distribuição dos retornos}

A Figura~\ref{fig:dist} compara o histograma dos retornos com a densidade da distribuição normal de mesma média e variância e apresenta o QQ-plot contra a normal. Os dois gráficos rejeitam, na prática, a hipótese de normalidade. O pico empírico em torno de zero é bem mais alto que o previsto pela normal teórica, e há massa de probabilidade nas caudas que a curva vermelha não consegue acomodar. No QQ-plot, os pontos extremos se afastam visivelmente da reta de referência, sobretudo na cauda inferior, onde aparecem retornos abaixo de $-30\%$ que estariam praticamente excluídos sob normalidade. Esses padrões reforçam a necessidade de modelar as inovações por uma distribuição com caudas pesadas.

\begin{figure}[H]
\centering
\includegraphics[width=\linewidth]{02_distribuicao.png}
\caption{Histograma com densidade normal teórica sobreposta (esquerda) e QQ-plot contra a normal (direita).}
\label{fig:dist}
\end{figure}

\subsection{Testes de hipótese}

A Tabela~\ref{tab:testes} consolida os testes estatísticos aplicados. O teste de Jarque-Bera rejeita fortemente a normalidade, com estatística superior a 30 mil. Os testes ADF e KPSS, aplicados de forma confirmatória, indicam estacionariedade dos log-retornos. O ADF rejeita a presença de raiz unitária com $p$-valor virtualmente nulo, enquanto o KPSS não rejeita a hipótese nula de estacionariedade ao nível de 5\%, com $p$-valor de $0{,}0783$ --- valor próximo do limite, o que merece atenção e foi um dos motivos para considerar quebras estruturais via análise de intervenção. O teste ARCH-LM com 10 defasagens rejeita de forma esmagadora a homocedasticidade, com $p$-valor da ordem de $10^{-136}$, justificando a adoção de modelos da família GARCH.

\begin{table}[H]
\centering
\caption{Testes estatísticos aplicados aos log-retornos.}
\label{tab:testes}
\begin{tabular}{lrrl}
\toprule
Teste & Estatística & $p$-valor & Conclusão \\
\midrule
Jarque-Bera          & $30589{,}85$ & $\approx 0$ & Rejeita normalidade \\
ADF (raiz unitária)  & $-22{,}1476$ & $\approx 0$ & Série estacionária \\
KPSS                 & $0{,}3973$   & $0{,}0783$  & Não rejeita estacionariedade \\
ARCH-LM(10)          & $663{,}42$   & $\approx 0$ & Forte heterocedasticidade \\
\bottomrule
\end{tabular}
\end{table}

\subsection{Funções de autocorrelação}

A Figura~\ref{fig:acf} apresenta as ACFs dos retornos e dos retornos ao quadrado. A primeira é praticamente compatível com ruído branco, com defasagens isoladas atingindo as bandas de confiança, sem padrão sistemático. A segunda exibe autocorrelações persistentes e significativas até defasagens elevadas. Esse contraste é a assinatura clássica do efeito ARCH e dá suporte adicional à modelagem da variância condicional.

\begin{figure}[H]
\centering
\includegraphics[width=0.9\linewidth]{03_acf.png}
\caption{ACF dos log-retornos (acima) e dos log-retornos ao quadrado (abaixo).}
\label{fig:acf}
\end{figure}

% ============================================================
% 5. MODELAGEM ARMA-GARCH
% ============================================================
\section{Construção do Modelo ARMA-GARCH}
\label{sec:modelo}

\subsection{Modelos candidatos e seleção}

Estimei, por máxima verossimilhança e com inovação $t$-Student padronizada, três modelos. O primeiro, M1, é um AR(1)-GARCH(1,1)-$t$, especificação básica que serve de referência. O segundo, M2, é um AR(1)-EGARCH(1,1)-$t$, que adiciona termo de assimetria à dinâmica da variância. O terceiro, M3, é um ARX(1)-EGARCH(1,1)-$t$ que inclui na média os três regressores de intervenção descritos na Seção~\ref{sec:metodo}.

A primeira tentativa, ainda na fase exploratória, foi um AR(1)-GARCH(1,1) com inovação normal. Os diagnósticos rejeitaram fortemente a normalidade dos resíduos padronizados e isso me levou a adotar a $t$-Student como distribuição padrão para todos os modelos comparados aqui. A Tabela~\ref{tab:comparacao} compara os três modelos finais por log-verossimilhança, AIC e BIC.

\begin{table}[H]
\centering
\caption{Comparação dos modelos estimados.}
\label{tab:comparacao}
\begin{tabular}{lrrr}
\toprule
Modelo & Log-Lik & AIC & BIC \\
\midrule
GARCH(1,1)-$t$              & $-9058{,}18$ & $18128{,}37$ & $18166{,}21$ \\
EGARCH(1,1)-$t$             & $-9040{,}92$ & $18095{,}84$ & $18139{,}99$ \\
EGARCH(1,1)-$t$ + Interv.   & $-9036{,}17$ & $18092{,}34$ & $18155{,}41$ \\
\bottomrule
\end{tabular}
\end{table}

O modelo M3 teve o menor AIC e a maior log-verossimilhança --- foi o escolhido como modelo final. A redução de AIC do M1 para o M2, de $32{,}5$ unidades, sugere fortemente efeito alavancagem na série, e o ganho adicional do M2 para o M3, de $3{,}5$ unidades, mostra a relevância das intervenções. O BIC, mais parcimonioso, classifica o EGARCH simples ligeiramente à frente, mas optei pelo M3 por causa da elevada significância estatística dos coeficientes de intervenção, apresentada a seguir.

\subsection{Estimativas do modelo final}

A Tabela~\ref{tab:coefs} traz os coeficientes estimados, erros-padrão robustos, estatísticas $t$ e $p$-valores do modelo final.

\begin{table}[H]
\centering
\caption{Coeficientes estimados do modelo final EGARCH(1,1)-$t$ com intervenções.}
\label{tab:coefs}
\small
\begin{tabular}{lrrrr}
\toprule
Parâmetro & Coef.\ & EP & $t$ & $p$-valor \\
\midrule
\multicolumn{5}{l}{\textit{Equação da média}} \\
Constante                       & $0{,}0456$  & $0{,}0483$  & $0{,}945$   & $0{,}345$ \\
AR(1)                           & $-0{,}0316$ & $0{,}0250$  & $-1{,}261$  & $0{,}207$ \\
$I_{\text{Castello}}$ (pulso)   & $-6{,}4969$ & $0{,}665$   & $-9{,}776$  & $1{,}42 \times 10^{-22}$ \\
$I_{\text{Lula}}$ (degrau)      & $0{,}0715$  & $0{,}0654$  & $1{,}092$   & $0{,}275$ \\
$I_{\text{Prates}}$ (pulso)     & $-2{,}0602$ & $0{,}221$   & $-9{,}312$  & $1{,}25 \times 10^{-20}$ \\
\midrule
\multicolumn{5}{l}{\textit{Equação da volatilidade}} \\
$\omega$                        & $0{,}0282$  & $0{,}00770$ & $3{,}664$   & $2{,}48 \times 10^{-4}$ \\
$\alpha$                        & $0{,}1389$  & $0{,}0188$  & $7{,}381$   & $1{,}58 \times 10^{-13}$ \\
$\gamma$                        & $-0{,}0301$ & $0{,}00900$ & $-3{,}340$  & $8{,}37 \times 10^{-4}$ \\
$\beta$                         & $0{,}9888$  & $0{,}00365$ & $271{,}06$  & $< 10^{-300}$ \\
\midrule
\multicolumn{5}{l}{\textit{Distribuição}} \\
$\nu$                           & $4{,}9073$  & $0{,}406$   & $12{,}092$  & $1{,}16 \times 10^{-33}$ \\
\bottomrule
\end{tabular}
\end{table}

Vale olhar os parâmetros com calma. O $\beta = 0{,}989$ mostra volatilidade muito persistente: choques na variância demoram a se dissipar, o que é típico de ações alavancadas em mercados emergentes. O $\alpha = 0{,}139$, significativo, mede a sensibilidade da variância aos choques recentes em valor absoluto. O $\gamma = -0{,}030$, com $p$-valor inferior a $0{,}001$, confirma o efeito alavancagem antes inferido pela queda de AIC: choques negativos elevam $\ln \sigma_t^2$ em maior magnitude que choques positivos de mesma intensidade, o que concorda com a literatura empírica em finanças \cite{morettin2017}. Por fim, a estimativa de $\nu = 4{,}91$ graus de liberdade é bastante baixa e indica caudas extremamente pesadas. Note-se que para $\nu < 5$ os momentos de ordem superior à quarta nem existem na distribuição teórica.

% ============================================================
% 6. ANÁLISE DE INTERVENÇÃO
% ============================================================
\section{Análise de Intervenção}
\label{sec:interv}

A demissão de Roberto Castello Branco, anunciada em 19/02/2021, gerou o maior efeito de intervenção identificado. O coeficiente do pulso correspondente foi estimado em $-6{,}50$ pontos percentuais, com $p$-valor da ordem de $10^{-22}$, indicando queda atípica do log-retorno do dia mesmo após controlar por toda a estrutura ARMA-EGARCH. O resultado reflete a percepção do mercado quanto ao risco de interferência política sobre a empresa --- percepção que se traduziu em forte revisão imediata dos preços.

Já a mudança da política de preços de combustíveis, anunciada em 08/03/2023 sob o governo Lula, foi modelada por dummy do tipo degrau. O coeficiente estimado, de $0{,}072$ ponto percentual ao dia, não foi significativo ($p = 0{,}275$). Isso sugere que, controlando-se pelos demais fatores do modelo, não houve alteração persistente do nível médio dos retornos após a mudança formal da política. Uma interpretação plausível é que parte relevante do impacto havia sido absorvida nos meses anteriores, à medida que a nova diretriz foi sendo sinalizada --- comportamento compatível com a hipótese de mercados eficientes na forma semiforte.

A demissão de Jean Paul Prates, em 14/05/2024, também resultou em coeficiente altamente significativo, com queda estimada de $2{,}06$ pontos percentuais no log-retorno do dia ($p \approx 10^{-20}$), porém de magnitude menor que o evento de 2021. A diferença pode ser atribuída à maior antecipação do mercado, dado que rumores sobre a saída do executivo precediam o anúncio formal.

% ============================================================
% 7. DIAGNÓSTICO
% ============================================================
\section{Diagnóstico do Modelo}
\label{sec:diag}

\subsection{Análise gráfica dos resíduos}

A Figura~\ref{fig:diag} traz o painel de diagnóstico dos resíduos padronizados do modelo final. O gráfico de série temporal dos resíduos não revela padrões sistemáticos remanescentes, e suas funções de autocorrelação --- tanto sobre os próprios resíduos quanto sobre seus quadrados --- ficam, em geral, dentro das bandas de confiança. O QQ-plot contra a $t$-Student com $\nu = 4{,}91$ apresenta bom alinhamento, exceto em poucos pontos extremos da cauda inferior, principalmente associados à pandemia da COVID-19. No conjunto, os gráficos indicam que o modelo se ajusta bem aos dados.

\begin{figure}[H]
\centering
\includegraphics[width=\linewidth]{04_diagnostico.png}
\caption{Painel de diagnóstico do modelo final: resíduos padronizados, suas ACFs e QQ-plot contra a $t$-Student.}
\label{fig:diag}
\end{figure}

\subsection{Volatilidade condicional}

A Figura~\ref{fig:vol} exibe a volatilidade condicional estimada. Os picos coincidem com episódios bem documentados: o ciclo da Operação Lava Jato e do choque do petróleo entre 2014 e 2016, a pandemia da COVID-19 em março de 2020, a turbulência associada à demissão de Castello Branco em fevereiro de 2021 e os movimentos relacionados às mudanças de comando e política regulatória da estatal nos anos posteriores. A persistência elevada captada pelo coeficiente $\beta$ é visível na lentidão com que as elevações de volatilidade voltam aos níveis médios.

\begin{figure}[H]
\centering
\includegraphics[width=\linewidth]{05_volatilidade.png}
\caption{Volatilidade condicional estimada pelo modelo EGARCH(1,1)-$t$ com intervenções.}
\label{fig:vol}
\end{figure}

% ============================================================
% 8. VAR
% ============================================================
\section{Estimação do Value-at-Risk e Backtesting}
\label{sec:var}

A partir do modelo final, calculei os limites diários de Value-at-Risk nos níveis de 95\% e 99\%, segundo

$$
\mathrm{VaR}_t(\alpha) = -\big( \hat{\mu}_t + q_\alpha(\nu)\, \hat{\sigma}_t \big),
$$

onde $q_\alpha(\nu)$ é o quantil-$\alpha$ da $t$-Student padronizada com $\nu = 4{,}91$ graus de liberdade. A Tabela~\ref{tab:var} resume o resultado do backtesting nos dois níveis.

\begin{table}[H]
\centering
\caption{Resultados do backtesting do VaR.}
\label{tab:var}
\begin{tabular}{lrrrrr}
\toprule
Nível & Violações & Taxa obs.\ & Taxa esp.\ & Kupiec LR & $p$-valor \\
\midrule
95\% & 197 & $4{,}86\%$ & $5{,}00\%$ & $0{,}170$ & $0{,}680$ \\
99\% &  45 & $1{,}11\%$ & $1{,}00\%$ & $0{,}481$ & $0{,}488$ \\
\bottomrule
\end{tabular}
\end{table}

No nível de 95\%, a taxa observada de violações ($4{,}86\%$) está muito próxima da nominal, com desvio inferior a 0,15 ponto percentual em uma amostra de 4053 dias, e o teste de Kupiec não rejeita cobertura correta. No nível de 99\%, o teste também não rejeita ($\mathrm{LR} = 0{,}481$, $p = 0{,}488$), o que valida o modelo para gestão de risco de mercado em padrões compatíveis com os Acordos de Basileia. A Figura~\ref{fig:var} mostra os limites de $-\mathrm{VaR}$ sobrepostos aos retornos observados, com as violações do nível de 99\% destacadas em vermelho.

\begin{figure}[H]
\centering
\includegraphics[width=\linewidth]{06_var.png}
\caption{Retornos diários, limites $-\mathrm{VaR}$ a 95\% e 99\% e violações do VaR 99\%.}
\label{fig:var}
\end{figure}

% ============================================================
% 9. DISCUSSÃO CRÍTICA
% ============================================================
\section{Discussão Crítica}
\label{sec:disc}

O modelo final tem virtudes importantes do ponto de vista estatístico e prático. O EGARCH com inovações $t$-Student capturou ao mesmo tempo clusters de volatilidade, leptocurtose e efeito alavancagem --- os principais fatos estilizados das séries financeiras. Todos os parâmetros têm interpretação econômica clara, o que dá transparência ao modelo e o torna defensável em contextos regulatórios. A inferência se baseou em erros-padrão robustos, e o backtesting empregou um teste com distribuição assintótica conhecida, permitindo decisões formais sobre adequação. O resultado positivo no teste de Kupiec a 99\% reforça a aplicabilidade prática.

Há limitações que vale comentar. O EGARCH suaviza a representação de saltos discretos de grande magnitude --- e a amostra tem um caso emblemático: o pico da pandemia, com retorno mínimo de $-35{,}24\%$. Modelos com componente de saltos, como o Jump-GARCH ou variantes com inovações de mistura, capturariam melhor esses extremos.

O coeficiente $\beta = 0{,}989$ está colado na fronteira de não-estacionariedade da variância, o que pode indicar quebras estruturais não totalmente capturadas pelas três intervenções modeladas. Modelos com mudança de regime markoviana, como o Markov-Switching GARCH, seriam uma alternativa direta \cite{morettin2020}.

Além disso, a constante e o coeficiente AR(1) na equação da média não foram significativos. O resultado é coerente com a hipótese de mercados eficientes na forma fraca, mas indica baixo poder preditivo do componente de média --- algo importante de explicitar para quem fosse usar o modelo para previsão e não apenas para risco. Por fim, o teste de Kupiec verifica apenas cobertura incondicional do VaR; um exame complementar de independência das violações ao longo do tempo (teste de Christoffersen) seria recomendável e fica como melhoria evidente para uma versão futura.

A coleta de dados também tem três pontos sensíveis. Primeiro, os preços de fechamento ignoram a dinâmica intraday, na qual frequentemente se concentram os picos de volatilidade durante anúncios e divulgações de resultados --- medidas baseadas em volatilidade realizada com dados de alta frequência seriam preferíveis. Segundo, ainda que o Yahoo Finance ajuste os preços retroativamente para proventos e desdobramentos, podem ficar descontinuidades em torno de eventos relevantes, sobretudo em datas ex-dividendo de grandes pagamentos. Terceiro, a presença dos retornos extremos da pandemia tende a dominar a estimação dos parâmetros e poderia ser tratada de forma específica, seja por marcação como observações influentes, seja por modelagem de regime distinto.

Como caminhos para extensões futuras, destaco: incluir regressores macroeconômicos exógenos --- preço do petróleo Brent, taxa de câmbio, taxa Selic --- na equação da média ou da variância; modelagem multivariada com ações correlacionadas (VALE3, ITUB4) via DCC-GARCH \cite{morettin2020}; teoria de valores extremos para estimar a cauda do VaR em níveis mais conservadores como $99{,}5\%$ ou $99{,}9\%$; modelos de volatilidade realizada com dados intraday; e estimação em janelas móveis para avaliar a estabilidade dos parâmetros ao longo do período amostral.

% ============================================================
% 10. CONCLUSÃO
% ============================================================
\section{Conclusão}
\label{sec:conc}

Apliquei a metodologia ARMA-GARCH à modelagem de retornos da PETR4 entre 2010 e 2026. Os resultados confirmam todos os fatos estilizados clássicos das séries financeiras: estacionariedade dos retornos, leptocurtose extrema, autocorrelação dos retornos ao quadrado e efeito alavancagem. O modelo final, EGARCH(1,1)-$t$ com regressores de intervenção, passou nos diagnósticos formais e no backtesting de VaR a 99\% pelo teste de Kupiec, com aderência adequada aos dados e potencial uso para gestão de risco de mercado.

A análise de intervenção mostrou que dois dos três eventos modelados --- as demissões de CEOs em 2021 e 2024 --- tiveram impactos negativos altamente significativos nos retornos do dia do anúncio. Já a mudança de política de preços de março de 2023 não apresentou efeito permanente estatisticamente detectável, resultado consistente com a hipótese de antecipação do mercado.

As limitações apontadas e as melhorias sugeridas, sobretudo a incorporação de variáveis exógenas e a modelagem multivariada, são caminhos diretos para a continuidade desta linha de pesquisa.

% ============================================================
% REFERÊNCIAS
% ============================================================
\begin{thebibliography}{9}

\bibitem{morettin2017}
MORETTIN, P.~A. \textit{Econometria Financeira: Um Curso em Séries Temporais Financeiras}. 3.\ ed.\ São Paulo: Edgard Blücher, 2017. ISBN: 978-85-212-1233-1.

\bibitem{morettin2018}
MORETTIN, P.~A.; TOLOI, C.~M.~C. \textit{Análise de Séries Temporais, Volume 1: Modelos Lineares Univariados}. 3.\ ed.\ São Paulo: Edgard Blücher (ABE -- Projeto Fisher), 2018. ISBN: 978-85-212-1888-3.

\bibitem{morettin2020}
MORETTIN, P.~A.; TOLOI, C.~M.~C. \textit{Análise de Séries Temporais, Volume 2: Modelos Multivariados e Não-Lineares}. São Paulo: Edgard Blücher (ABE -- Projeto Fisher), 2020. ISBN: 978-65-5506-029-0.

\bibitem{shumway2017}
SHUMWAY, R.~H.; STOFFER, D.~S. \textit{Time Series Analysis and Its Applications: With R Examples}. 4.\ ed.\ New York: Springer, 2017. ISBN: 978-3-319-52451-1.

\bibitem{engle1982}
ENGLE, R.~F. Autoregressive conditional heteroscedasticity with estimates of the variance of United Kingdom inflation. \textit{Econometrica}, v.~50, n.~4, p.~987--1007, 1982. ISSN: 0012-9682.

\bibitem{bollerslev1986}
BOLLERSLEV, T. Generalized autoregressive conditional heteroskedasticity. \textit{Journal of Econometrics}, v.~31, n.~3, p.~307--327, 1986. ISSN: 0304-4076.

\bibitem{nelson1991}
NELSON, D.~B. Conditional heteroskedasticity in asset returns: a new approach. \textit{Econometrica}, v.~59, n.~2, p.~347--370, 1991. ISSN: 0012-9682.

\bibitem{sheppard2024}
SHEPPARD, K. \textit{arch: Autoregressive Conditional Heteroskedasticity (ARCH) and other tools for financial econometrics, written in Python}. Documentação do pacote, 2024. Disponível em: \url{https://bashtage.github.io/arch/}.

\end{thebibliography}

\end{document}
